\worksheettitle
  {Домашняя практика}
  {\modulemark{M04-H} Точки на отрезке}

\studentnameblock

\begin{practicebox}[1. Целое и части]
  Точка $K$ лежит между $M$ и $N$.
  \begin{enumerate}[label=\alph*)]
    \item $MK=4$ см $6$ мм, $KN=3$ см $8$ мм. Найдите $MN$.
    \item $MN=15$ см, $MK=8$ см $4$ мм. Найдите $KN$.
  \end{enumerate}
  \writinglines{3}
\end{practicebox}

\begin{practicebox}[2. Две внутренние точки]
  Порядок $A-C-D-B$. $AB=24$ см, $AC=9$ см, $DB=6$ см.
  Запишите равенство частей и найдите $CD$. \writinglines{3}
\end{practicebox}

\begin{practicebox}[3. Один общий конец]
  Порядок $P-R-S$. $PR=5$ см $7$ мм, $PS=11$ см $2$ мм.
  Найдите $RS$. \writinglines{2}
\end{practicebox}

\begin{warningbox}[4. Исправьте ошибку]
  При порядке $K-L-M$, $KL=7$ см и $KM=12$ см ученик получил
  $LM=19$ см. Назовите причину и запишите исправление.
  \writinglines{2}
\end{warningbox}
